% !TeX root = ../main.tex
% sections/01_introduction.tex -- Section I of the MobiWac 2026 camera-ready (accepted 2026-08-26).
% Base text: the dissertation's Chapter 5 (articles/dissertacao/src/chapters/5_mobiwac/01_introduction.tex),
% ported 2026-09-06 (CONSOLIDATION_PLAN 3.6). Camera-ready edits carry a dated comment beside them.
\section{Introduction}
\label{sec:intro}

Location-based social networks, such as Gowalla~\cite{cho2011gowalla}, let people share the places that they visit as check-ins, which together record how people move through a city. Individual mobility is highly predictable in principle~\cite{song2010limits}, so a service that anticipates the next move can prepare ahead of time, caching content where a user is heading~\cite{bastug2014edge} or provisioning capacity at the destination before demand arrives. Handover predictions already let cellular services adapt in advance at the network level~\cite{vielhaus2022handover}; at the city scale, check-in mobility analysis is likewise a design input for mobility-aware services~\cite{moura2025mobilityaware}.

Predicting the next place exactly, a single point of interest (POI), is hard and often more than a service needs. Two coarser questions are usually enough: the next category, the type of place such as food or shopping, and the next region, the part of the city that a user is heading to (a neighborhood-scale unit such as the U.S. census tract). The first captures intent, the second captures geography; we study whether one model should learn both at once.

Sharing one representation across tasks has a cost: in multi-task learning (MTL), one model does several jobs at once by sharing most of its parts, so the shared parameters can converge to a compromise optimal for neither task, helping one while hurting the other~\cite{caruana1997multitask}. Our earlier
work reported no consistent multi-task advantage for the paired category tasks and attributed it, in part, to this effect~\cite{silva2025mtlnet}; the useful question is where sharing helps, what it costs, and how to share so the gains hold and the cost stays small.
% [B.1 FIX 2026-07-24, author-approved] Was "Prior work observed exactly this for next-category
% and next-region": two errors vs the CBIC record (Ch.3) -- CBIC paired STATIC category
% classification with next-category (no region task), and it HYPOTHESIZED negative transfer on a
% parity null, it did not OBSERVE it. Repaired per persona-14 B.1. Mirrored in [mobiwac]/src/ and
% logged in [mobiwac]/ERRATA.md + Appendix B.

% [2026-09-02, B-23, author ruling "concordo podemos fazer"] "We propose two enhancements." named no
% antecedent; the previous paragraph's \cite{silva2025mtlnet} is it. Reproduced accepted-article text
% -> errata row B.5 (registry).
We propose two enhancements to that earlier model. First, we build a check-in-level representation: instead of one fixed vector per place, each check-in gets its own vector that holds its context (the time, nearby places, and recent visits).
This adds a fourth level, the check-in, beneath the place, region, and city levels of hierarchical graph infomax~\cite{huang2023hgi,velickovic2019dgi} (Fig.~\ref{fig:dataflow}). Second, we train one model that
predicts the next category and the next region in a single forward pass, with a shared trunk (a cross-attention stack where the two tasks exchange
semantic context) and a private spatial path for the region task (Fig.~\ref{fig:model}). We evaluate on two settings chosen to differ: five U.S. states of different sizes (Gowalla) and one non-U.S. city (Istanbul).\footnote{\label{fn:code}Code (model, representation, baselines, statistical tests): \url{https://github.com/VitorHugoOli/PoiMtlNet/tree/mobiwac}. Both data sources are public: the category-annotated Gowalla dump (\url{https://figshare.com/articles/dataset/gowalla_data/22126586}) and Massive-STEPS~\cite{wongso2025massivesteps}.}

Our contributions are as follows.
\begin{itemize}
  \item \textbf{A check-in-level representation.} We extend hierarchical graph infomax from the
    place to the individual visit, so that each check-in carries its own vector rather than sharing
    one fixed vector per place. Under a controlled comparison in which only the input changes, this
    improves next-category prediction at every one of the six datasets and in every one of their
    folds, by $+0.23$ to $+6.29$ points of macro-F1, and a paired test separates the two
    representations at five of the six (Table~\ref{tab:substrate}). A single vector per place cannot
    separate places that serve several purposes, and per-visit context recovers that distinction.
    The complete system is also above every external baseline reported, on both tasks and at every
    dataset, three of them re-run under our protocol (Section~\ref{sec:results-external}).
  \item \textbf{A single model for both tasks.} One model predicts the next category and the next
    region in one forward pass, sharing semantic context while keeping a private spatial path for
    the region task. To our knowledge, this is the first work to treat fine-grained region as an end
    target of equal standing (Section~\ref{sec:related-tasks}). It outperforms a dedicated
    single-task region model at the two datasets with the largest region vocabularies and stays
    within a registered two-point margin at the other four. At one of the two, a dedicated model of
    the same size scores above it, so we do not attribute that gain to sharing
    (Section~\ref{sec:discussion}). On category it outperforms the dedicated model at one dataset,
    and the five remaining differences are equivalent to zero within half a point, even though the
    dedicated category model receives a per-dataset search (Table~\ref{tab:results}).
  \item \textbf{Next-region results under one protocol.} We report next-region prediction across
    six datasets, four random initializations, and five folds, against three external systems and
    a Markov-1 floor, under a pre-registered statistical protocol, so that others can report on the
    same footing. The two datasets where the joint model outperforms are the two with the largest
    region vocabularies, while the four smaller ones sit inside the margin (Table~\ref{tab:results}).
    We report this as an observation and not a law, since the ordering does not hold inside the
    pair and region count co-varies with corpus size.
\end{itemize}
% [2026-09-06, camera-ready, REBALANCE, author-approved: PROPOSTA_REEQUILIBRIO.md section B]
% Three bullets in the order of the strength of the evidence. Departures from the proposal text:
% bullet 1's external clause said "the four external baselines re-run under our protocol"; ReHDM is
% reported under its own published protocol (Section V-D), so the clause now says "every external
% baseline reported ... three of them re-run under our protocol" (POI-RGNN, HMT-GRN, STAN); bullet
% 3's "against four external baselines" -> "three external systems and a Markov-1 floor" (the
% region column holds HMT-GRN, STAN, ReHDM and the floor); "significant" -> "a paired test
% separates" (test named); "reaches that gain" -> "scores above it"; "where joint training pays"
% -> "where the joint model outperforms" (money metaphor, GLOSSARY section 8); semicolons split.
% WITHHELD by condition 1 of the proposal: the clause "with the folds and region assignments
% released" does not enter until origin/mobiwac carries the reported generation; add it to bullet 3
% after "protocol" once the branch is republished. The sentence on seeds per table (former bullet 3)
% moved out: Section V and the Table II footnote state it. 471 -> ~300 words.

The rest of the paper presents related work (Section~\ref{sec:related}), the problem (Section~\ref{sec:problem}), the method and setup (Sections~\ref{sec:method} and~\ref{sec:setup}), results (Section~\ref{sec:results}), and discussion and conclusions (Sections~\ref{sec:discussion} and~\ref{sec:conclusion}).
